\silentchapter{Preface}{preface}
\begin{widebody}
\begin{foreordbody}

\begin{center}\itshape
This book is a drastic measure. The aim is not to improve the discipline of design;\\ it is to establish the ground the discipline can stand on.
\end{center}

Design lacks a grammar.

Why do things have the form they have? For more than a century the field has leaned on a variant of Sullivan's formula, \textit{form follows function}. Michl\textsuperscript{15} has shown that the formula explained everything and therefore nothing. Semper\textsuperscript{1} linked material technique, form, and cultural meaning, but the discipline inherited the split rather than the synthesis. Thompson\textsuperscript{2} showed that biological form follows physical forces; design never took the step from biology to artefact. The traditions have described the world of forms piecewise; none has unified them.

Without a common grammar the field has no internal way of saying no to what should not be made. The profession has no spine, neither collective nor academic. Targeting civilians, and interfaces engineered to extract attention from children, are design projects in the technical sense, and the field has no language to say anything else. This is not a moral failing in the individual designer; it is a structural hole in the discipline. A discipline without a foundation will not survive the coming decades. It will become obsolete, as phrenology did before neuroscience, as alchemy did before chemistry, as vitalism did before biology.

The framework can be summarised in one sentence: \textbf{\textit{form is a position in a space of possibilities, shaped by simultaneous selection pressures, navigated by agents at several scales.}}

Hence three layers. \textit{Form} is the ontological primitive. \textit{Form-maker} is the agent that navigates, substrate-independent: cell, tree, craftsperson. \textit{Designer} is the specific kind of form-maker who navigates on behalf of others who will live with the form.

The framework rests on three pillars. The \textbf{shape grammar} of Stiny and Gips\textsuperscript{11} supplies the generative layer. The \textbf{concept--knowledge theory} of Hatchuel and Weil\textsuperscript{16,\,17} formalises the epistemic navigation. The substrate-independent \textbf{agency framework} of Levin and Fields\textsuperscript{21,\,22} provides the ontological ground. That three independent traditions converge on the same structure confirms that the definition grasps something real.\hfill\textit{Oslo, 2026}
\end{foreordbody}
\end{widebody}
